% Lentille mince convergente : foyers, rayons particuliers et image réelle (L'essentiel, chapitre 18)
% f' = 1,5 ; OA = -4,0 ; AB = 1,2 ; OA' = 2,4 ; A'B' = -0,72
\begin{tikzpicture}[x=1.5cm,y=1.5cm,
  rayon/.style={otOrange!85!black,line width=0.8pt,postaction={decorate},
    decoration={markings,mark=at position 0.3 with {\arrow{Stealth[length=1.6mm]}}}}]
\draw[otInk!60,line width=0.5pt] (-4.8,0) -- (3.6,0);
\draw[otNavy,line width=1.2pt,{Stealth[length=2.2mm]}-{Stealth[length=2.2mm]}] (0,-1.35) -- (0,1.45);
\foreach \x/\n in {-1.5/F,1.5/F'}{\draw[otInk,line width=0.6pt] (\x,-0.06) -- (\x,0.06); \node[pcLegende,anchor=north] at (\x,-0.05) {$\n$};}
\node[pcLegende,anchor=north east] at (0,-0.02) {O};
\draw[otNavy,line width=1.2pt,-{Stealth[length=2.2mm]}] (-4,0) -- (-4,1.2) node[pcLegende,anchor=south,text=otNavy] {B};
\node[pcLegende,anchor=north] at (-4,-0.03) {A};
\draw[otGreenDark,line width=1.2pt,-{Stealth[length=2.2mm]}] (2.4,0) -- (2.4,-0.72) node[pcLegende,anchor=north,text=otGreenDark] {B'};
\node[pcLegende,anchor=south] at (2.4,0.03) {A'};
\draw[rayon] (-4,1.2) -- (0,1.2) -- (3.3,-1.44);
\draw[rayon] (-4,1.2) -- (3.3,-0.99);
\draw[rayon] (-4,1.2) -- (0,-0.72) -- (3.3,-0.72);
\draw[otInk!60,line width=0.4pt,{Stealth[length=1.4mm]}-{Stealth[length=1.4mm]}] (0,-1.15) -- (1.5,-1.15) node[pos=0.5,pcLegende,anchor=north] {$f'$};
\node[pcLegende,anchor=west,align=left] at (3.65,1.0) {1. parallèle à l'axe :\\ressort par $F'$\\[1mm]2. par O : non dévié\\[1mm]3. par $F$ : ressort\\parallèle à l'axe};
\end{tikzpicture}
\par\smallskip\legende{L'image $\mathrm{A'B'}$ d'un objet placé avant le foyer $F$ est réelle et renversée ; on peut la recueillir sur un écran.}
